\documentclass[11pt]{article}

\usepackage[letterpaper,margin=1in]{geometry}
\usepackage{amsmath,amssymb}
\usepackage{booktabs}
\usepackage{longtable}
\usepackage{array}
\usepackage{xcolor}
\usepackage{listings}
\usepackage{url}
\usepackage{hyperref}
\usepackage{titlesec}
\usepackage{parskip}
\usepackage{fancyhdr}
\usepackage{enumitem}

\hypersetup{
  colorlinks=true,
  linkcolor=black,
  urlcolor=blue!50!black,
  citecolor=black,
  pdftitle={KV Cache Engine --- Software Overview},
  pdfauthor={LonghornSilicon},
  pdfsubject={Compiler Co-Design Orientation}
}

\titleformat{\section}{\normalfont\Large\bfseries}{\thesection.}{0.6em}{}
\titleformat{\subsection}{\normalfont\large\bfseries}{\thesubsection}{0.6em}{}

\pagestyle{fancy}
\fancyhf{}
\lhead{\small KV Cache Engine --- Software Overview}
\rhead{\small v0.1}
\cfoot{\thepage}
\renewcommand{\headrulewidth}{0.4pt}

\definecolor{codebg}{gray}{0.97}
\definecolor{codekey}{rgb}{0.13,0.29,0.53}
\definecolor{codecom}{rgb}{0.25,0.5,0.25}
\lstset{
  basicstyle=\ttfamily\small,
  backgroundcolor=\color{codebg},
  frame=single,
  framerule=0.4pt,
  rulecolor=\color{black!30},
  xleftmargin=0.5em,
  xrightmargin=0.5em,
  aboveskip=0.6em,
  belowskip=0.6em,
  showstringspaces=false,
  breaklines=true,
  keywordstyle=\color{codekey}\bfseries,
  commentstyle=\color{codecom}\itshape,
  numbers=none
}

\setlist{itemsep=1pt,topsep=2pt}

\begin{document}

%==============================================================================
\thispagestyle{empty}
\begin{center}
\vspace*{1.2in}
{\Huge\bfseries KV Cache Engine}\\[0.4em]
{\LARGE\bfseries Software Overview}\\[1.2em]

{\large \textsc{LonghornSilicon --- Compiler Co-Design Orientation}}\\[0.4em]
{\large Block 2 of 4 --- KV Cache Compression / Decompression}\\[2em]

\rule{0.7\textwidth}{0.4pt}\\[1em]

\begin{tabular}{rl}
\textbf{Version}:    & 0.1 \\
\textbf{Date}:       & 2026-05-14 \\
\textbf{Audience}:   & Compiler team, runtime engineers \\
\end{tabular}\\[1em]
\rule{0.7\textwidth}{0.4pt}
\vfill
\end{center}
\clearpage

%==============================================================================
\tableofcontents
\clearpage

%==============================================================================
\section{Purpose}

This document orients the compiler team on the KV cache engine block.
It covers: what the block does, how to use the reference models, what
the compiler should emit, and what integration phases look like.

For the formal interface contract, see the ISA specification
(\path{docs/isa/kv_cache_engine_isa.pdf}). For the API reference, see
\path{docs/reference_model_api.pdf}.

%==============================================================================
\section{What the Block Does}

The KV cache engine sits between the attention compute unit (block~1)
and the memory hierarchy controller (block~4). It compresses key and
value vectors before writing to on-chip SRAM, and decompresses them
when the attention unit needs them for dot-product computation.

\begin{itemize}
\item \textbf{On write}: incoming FP16 K/V vectors are compressed via
      TurboQuant+ (PolarQuant + QJL for keys, PolarQuant only for values)
      and stored in SRAM.
\item \textbf{On read}: compressed entries are decompressed back to
      FP16 coordinates and streamed to the attention unit via AXI-Stream.
\item \textbf{Eviction}: when SRAM is full, the block signals the memory
      hierarchy controller to spill entries to DRAM.
\end{itemize}

\subsection{Compression Pipeline}

\begin{enumerate}
\item \textbf{Norm extraction} --- L2 norm stored as metadata (16-bit fixed-point)
\item \textbf{Rotation} --- Walsh-Hadamard Transform + random sign flips
      (deterministic from seed, self-inverse, $O(n \log n)$)
\item \textbf{Quantization} --- 3-bit optimal Lloyd-Max scalar quantization
      (8 centroids symmetric around zero)
\item \textbf{QJL projection} (K only) --- 1-bit sign projection of the
      quantization residual (preserves inner products)
\item \textbf{Packing} --- bit-pack norm + indices + residual norm + signs
      into a compact SRAM word
\end{enumerate}

\noindent Decompression reverses the pipeline: unpack $\rightarrow$
centroid lookup $\rightarrow$ QJL correction (K only) $\rightarrow$
inverse WHT $\rightarrow$ norm rescaling.

%==============================================================================
\section{Reference Models}

\subsection{Directory Layout}

\begin{lstlisting}[language=bash]
sw/
+-- README.md                     <- orientation (this doc's markdown form)
+-- reference_model/
    +-- README.md                 <- API reference + build instructions
    +-- Makefile
    +-- kv_cache_engine_ref.py    <- Python golden reference
    +-- kv_cache_engine_ref.hpp   <- C++ header
    +-- kv_cache_engine_ref.cpp   <- C++ implementation
    +-- test_kv_cache_engine_ref.py
    +-- test_kv_cache_engine_ref.cpp
\end{lstlisting}

\subsection{Which Models Are Available}

\begin{table}[h]
\centering
\small
\begin{tabular}{@{}l l l l l@{}}
\toprule
Block & Status & C++ class & C API & Python \\
\midrule
KV Cache Engine & Bit-exact vs RTL & \texttt{lhsi::KVCacheEngine} & \texttt{lhsi\_kv\_*} & \texttt{KVCacheEngine} \\
\bottomrule
\end{tabular}
\end{table}

\subsection{Build Everything}

\begin{lstlisting}[language=bash]
cd sw/reference_model
make test-all     # builds C++ tests, runs them, runs Python tests
make shared       # libkv_cache_engine_ref.so
make static       # libkv_cache_engine_ref.a
make clean
\end{lstlisting}

\noindent Requires Python 3.10+, NumPy, and a C++17 compiler.

%==============================================================================
\section{What the Compiler Should Emit}

A compiler backend targeting the KV cache engine emits sequences of
the operations defined in the ISA spec (\S4 of
\path{docs/isa/kv_cache_engine_isa.pdf}):

\begin{table}[h]
\centering
\small
\begin{tabular}{@{}l l p{6cm}@{}}
\toprule
Operation & C API & What it does \\
\midrule
\texttt{KV\_QUERY} & \texttt{lhsi\_kv\_info()} & Read hardware config \\
\texttt{KV\_RESET} & \texttt{lhsi\_kv\_reset()} & Clear SRAM, reset FSM \\
\texttt{KV\_COMPRESS\_KEY} & \texttt{lhsi\_kv\_compress\_key()} & Compress + store key \\
\texttt{KV\_COMPRESS\_VALUE} & \texttt{lhsi\_kv\_compress\_value()} & Compress + store value \\
\texttt{KV\_DECOMPRESS\_KEY} & \texttt{lhsi\_kv\_decompress\_key()} & Read + decompress key \\
\texttt{KV\_DECOMPRESS\_VALUE} & \texttt{lhsi\_kv\_decompress\_value()} & Read + decompress value \\
\bottomrule
\end{tabular}
\end{table}

\subsection{Worked Example: Attention with KV Cache}

\begin{lstlisting}[language=Python]
from kv_cache_engine_ref import KVCacheEngine

engine = KVCacheEngine()

# Store phase: compress new tokens' K/V
for token_idx, (k, v) in enumerate(new_kv_pairs):
    addr = token_idx % engine.info.sram_depth
    engine.store_kv(addr, k, v)

# Attention phase: decompress for Q*K^T
for cached_idx in range(num_cached):
    k_hat = engine.load_k(cached_idx)
    v_hat = engine.load_v(cached_idx)
    score = dot(query, k_hat) / sqrt(D)
    # Route through precision controller (block 1)
    # ...
\end{lstlisting}

%==============================================================================
\section{Integration Phases}

\begin{table}[h]
\centering
\small
\renewcommand{\arraystretch}{1.2}
\begin{tabular}{@{}c p{3cm} p{4cm} p{4.5cm}@{}}
\toprule
Phase & Timeline & Compiler targets & We provide \\
\midrule
0 & now & Python + C++ reference models & Models + ISA + tests \\
1 & when FPGA arrives & AXI-Lite + AXI-Stream on Zynq & Vivado bitstream + driver \\
2 & all 4 blocks done & Multi-block FPGA project & Integrated attention pipeline \\
3 & post-tape-out & PCIe-attached chip & TSMC 16FFC silicon + driver \\
\bottomrule
\end{tabular}
\end{table}

\noindent The ISA contract (\path{docs/isa/kv_cache_engine_isa.pdf}) is
stable across all four phases. Only the runtime implementation changes ---
Python function call, AXI register write, or PCIe MMIO.

%==============================================================================
\section{Interaction with Block~1 (Precision Controller)}

The KV cache engine and precision controller work together during
attention inference:

\begin{enumerate}
\item KV cache engine decompresses K/V vectors from SRAM
\item The attention unit computes $Q \cdot K^T$ scores
\item The precision controller decides INT8 vs FP16 per tile
\item The MAC array computes $\mathrm{softmax}(S) \cdot V$ at the chosen precision
\end{enumerate}

\noindent A compiler backend that targets both blocks emits interleaved
\texttt{KV\_DECOMPRESS\_*} and \texttt{PC\_PUSH\_TILE} operations. The
reference models for both blocks can be composed in the same process:

\begin{lstlisting}[language=Python]
from kv_cache_engine_ref import KVCacheEngine
from precision_controller_ref import PrecisionController  # from block 1

kv = KVCacheEngine()
pc = PrecisionController()

k_hat = kv.load_k(addr)
scores = query @ k_hat  # Q * K^T
decision = pc.process_tile(scores)
# decision: True = FP16, False = INT8
\end{lstlisting}

%==============================================================================
\section{Open Questions for the Compiler Team}

\begin{enumerate}
\item \textbf{Granularity}: does the compiler emit individual compress/decompress
      ops, or a fused \texttt{kv\_attention(Q, K\_cache, V\_cache)} op?
\item \textbf{Batch interface}: should the reference model support batched
      compress/decompress (multiple vectors in one call) for throughput?
\item \textbf{Async model}: should \texttt{KV\_COMPRESS\_KEY} be non-blocking
      with a separate completion query, matching the hardware pipeline?
\item \textbf{Cost model}: are placeholder cycle counts sufficient, or do
      you need post-synthesis numbers for the scheduler?
\end{enumerate}

\vspace{0.5in}
\noindent\rule{\textwidth}{0.4pt}\\
\noindent\small\textit{LonghornSilicon KV Cache Engine --- Software Overview.
This document is open and may be redistributed.}

\end{document}
